\section{Entropy roofs and the whole weighted adjacency}
\label{sec:operator}

We now put computation and periodic support into one operator.  Freeze the
edge roof
\begin{align}
\tau(I_n,T_{n,2})&=\log(2n),\label{eq:roof-input}\\
\tau(T_{n,d},A_n)&=\log(nd),\label{eq:roof-accept}\\
\tau(T_{n,d},Q_{n,d,2})&=\log(2nd),\label{eq:roof-tq}\\
\tau(Q_{n,d,q},\mathrm{next})&=\log(ndq),\label{eq:roof-q}\\
\tau(R_{n,k},R_{n,k+1})&=\log(n(k+1)),\label{eq:roof-r}\\
\tau(A_n,A_n)&=\log n.\label{eq:roof-loop}
\end{align}
Every value is the entropy of a tensor product of source full shifts.  The
factor $2$ in \eqref{eq:roof-tq} records entry into the first explicit
quotient state and matches the finite implementation.

There are two different naturality claims.  The accepted loop roof
$h(F_n)=\log n$ is canonical on the chosen source skeleton, but it is
attached only after the program accepts $n$.  By contrast, the transient
roofs \eqref{eq:roof-input}--\eqref{eq:roof-r} are a
\texttt{MODELING\_CHOICE}.  They are source-expressible and useful for
summability, but trial division does not force them.  Any locally uniformly
summable replacement on transient edges gives the same power traces and
determinant.

Let
\[
\Hcal=\ell^2(V(G))
\]
with standard basis $(\delta_v)_{v\in V(G)}$.  For $\RePart s>1$, define the
arrival-weighted adjacency
\begin{equation}
L_s\delta_u=\sum_{e:u\to v}e^{-s\tau(e)}\delta_v.
\label{eq:operator}
\end{equation}
The sum has one term at each source vertex because the graph is
deterministic.  We retain the edge-sum notation to expose the nuclear
decomposition.  No thermodynamic function space is being used, so we call
$L_s$ a weighted vertex-adjacency or graph transfer, not a Ruelle operator.

\begin{theorem}[Trace class and holomorphic family]
\label{thm:traceclass}
For every $s$ with $\sigma=\RePart s>1$, the operator $L_s$ is trace class.
The map $s\mapsto L_s$ is holomorphic with values in $\Sone(\Hcal)$ on that
half-plane.
\end{theorem}

\begin{proof}
Each edge contributes the rank-one operator
\[
e^{-s\tau(e)}|\delta_{t(e)}\rangle\langle\delta_{o(e)}|
\]
of trace norm $e^{-\sigma\tau(e)}$.  The loop and input sums are bounded by
\[
\sum_p p^{-\sigma}+\sum_{n\ge2}(2n)^{-\sigma}<\infty.
\]
For divisor-state edges, including the unique terminal value
$d=\lfloor\sqrt n\rfloor+1$, a fixed endpoint constant gives
\[
\sum_{n\ge2}n^{-\sigma}
 \sum_{2\le d\le\sqrt n+1}d^{-\sigma}
\le C_\sigma\sum_{n\ge2}n^{-\sigma}<\infty.
\]
For quotient states, enlarging every finite reachable range yields
\[
\sum_{n\ge2}\sum_{d\ge2}\sum_{q\ge2}(ndq)^{-\sigma}
\le
\left(\sum_{n\ge2}n^{-\sigma}\right)
\left(\sum_{d\ge2}d^{-\sigma}\right)
\left(\sum_{q\ge2}q^{-\sigma}\right)<\infty.
\]
Finally the cemetery contribution is bounded by
\[
\sum_{\substack{n\ge2\\ n\,\mathrm{composite}}}
\sum_{k\ge1}[n(k+1)]^{-\sigma}
\le
\left(\sum_{n\ge2}n^{-\sigma}\right)
\left(\sum_{j\ge2}j^{-\sigma}\right)<\infty.
\]
Thus the edge series converges absolutely in trace norm.  Replacing
$\sigma$ by a compact lower bound on every compact sub-half-plane gives
locally uniform convergence of holomorphic finite partial sums.  The
Banach-valued Weierstrass theorem proves $\Sone$-holomorphy.
\end{proof}

The theorem concerns the raw graph, not a recurrent pruning.  This is the
strongest analytic success of the candidate: computation edges, cemetery
edges, and accepted loops genuinely coexist in one trace-class operator on
the same half-plane as the Euler product.
